\chapter{AC Steady State: Impedance, Admittance, and Power}
\label{ch:ac}
The previous chapters introduced the components---\(R\), \(L\), \(C\), and
sources---and \cref{ch:complex} introduced the phasor as the complex number that
stands in for a steady sinusoid. This chapter puts the two together. In the
\textbf{AC steady state}---a circuit driven by a sinusoid long enough for
transients to die out---every voltage and every current is a sinusoid at the
source frequency, so each can be written as a phasor. The element laws, rewritten
in phasors, then give the central circuit quantities: impedance, admittance, and
complex power. This is the phasor machinery of \cref{ch:complex} applied to actual
circuits.

\section{From Element Laws to Impedance}
Recall the payoff of \cref{ch:complex}: for a phasor quantity, \(\dd/\dd t\)
becomes multiplication by \(\jj\omega\). Apply that to each element law. The
resistor's \(v=Ri\) is already algebraic. The capacitor's \(i=C\,\dd v/\dd t\)
becomes \(I=\jj\omega C\,V\), and the inductor's \(v=L\,\dd i/\dd t\) becomes
\(V=\jj\omega L\,I\). In every case the phasor voltage is a complex multiple of
the phasor current---exactly the algebraic form of Ohm's law, with a complex
proportionality constant.

That constant is the \textbf{impedance}, the complex ratio of the voltage phasor
to the current phasor,
\[
Z=\frac{V}{I}=R+\jj X,
\]
a frequency-dependent generalization of resistance. Its real part \(R\) is
resistance (dissipation); its imaginary part \(X\) is \textbf{reactance} (energy
storage, no dissipation). For the three elements,
\[
Z_R=R,\qquad
Z_L=\jj\omega L,\qquad
Z_C=\frac{1}{\jj\omega C}=-\jj\frac{1}{\omega C}.
\]
An inductor has positive reactance \(X_L=\omega L\); a capacitor has negative
reactance \(X_C=1/(\omega C)\). At low frequency a capacitor's impedance is large
(it blocks DC) and an inductor's is small (it passes DC); at high frequency the
roles reverse.

\begin{physicalbox}
\textbf{A chemical-engineering analogy.} Circuits obey the same
across/through (effort/flow) structure as the hydraulic and thermal systems of
process dynamics:
\begin{center}
\begin{tabularx}{0.92\textwidth}{@{}L{0.26\textwidth}Y@{}}
\toprule
Electrical & Process / hydraulic analog \\
\midrule
Voltage \(V\) (across) & Pressure or head (driving force) \\
Current \(I\) (through) & Volumetric flow rate \\
Resistance \(R\) & Flow resistance (pressure drop per unit flow) \\
Capacitance \(C\) & Tank capacitance (volume stored per unit head) \\
Inductance \(L\) & Fluid inertance (inertia opposing flow change) \\
Impedance \(Z\) & Frequency-domain ``flow resistance'' \\
\bottomrule
\end{tabularx}
\end{center}
So \emph{impedance} is the generalized flow resistance of a network in the
frequency domain, and \emph{reactance} is the piece contributed by storage---a
surge tank or fluid inertia opposes a \emph{change} in flow without dissipating
energy, storing it and handing it back a quarter cycle later. That energy return,
not loss, is exactly why reactance shifts phase by \(90^\circ\) instead of
consuming power.
\end{physicalbox}

\section{Series and Parallel: Where the RLC Impedance Comes From}
Impedances in series add, exactly like resistors, because the same current flows
through each and the voltages sum (KVL). For \(R\), \(L\), \(C\) in series that
gives the impedance directly:
\[
Z=Z_R+Z_L+Z_C=R+\jj\omega L+\frac{1}{\jj\omega C}
=R+\jj\left(\omega L-\frac{1}{\omega C}\right).
\]
The reactance is the \emph{difference} \(\omega L-1/(\omega C)\); it vanishes at
resonance, a fact \cref{ch:rlc} builds on.

\begin{figure}[htbp]
\centering
\begin{circuitikz}
\draw (0,0) to[sV,l=\(v_s\)] (0,2)
      to[R,l=\(R\)] (2.6,2)
      to[L,l=\(L\)] (5.2,2)
      to[C,l=\(C\)] (5.2,0)
      -- (0,0);
\end{circuitikz}
\caption{A series RLC circuit whose impedance
\(Z=R+\jj(\omega L-1/\omega C)\) is just the sum of the three element
impedances.}
\end{figure}

\section{Admittance and Susceptance}
For parallel networks it is easier to work with the reciprocal of impedance, the
\textbf{admittance}
\[
Y=\frac{1}{Z}=G+\jj B,
\]
whose real part \(G\) is conductance and imaginary part \(B\) is
\textbf{susceptance}. Just as series impedances add, \emph{parallel admittances
add}, because the same voltage appears across each branch and the currents sum
(KCL):
\[
Y_{\mathrm{eq}}=\sum_k Y_k.
\]
In the process analogy, admittance is ``flow per unit driving force''---a
conductance to flow---and parallel pipes add their flow conductances the same way.

\subsection{Worked Example: a Parallel Branch}
A \(\SI{100}{\ohm}\) resistor is in parallel with a capacitor whose reactance is
\(\SI{100}{\ohm}\) at the frequency of interest. Their admittances are
\[
Y_R=\frac{1}{100}=\SI{0.010}{\siemens},
\qquad
Y_C=\frac{1}{-\jj100}=\jj\,\SI{0.010}{\siemens},
\]
so \(Y_{\mathrm{eq}}=0.010+\jj0.010=\SI{0.0141}{\siemens}\angle45^\circ\), and the
parallel impedance is
\[
Z_{\mathrm{eq}}=\frac{1}{Y_{\mathrm{eq}}}\approx\SI{70.7}{\ohm}\angle{-45^\circ}
=50-\jj50\ \Omega.
\]
Adding admittances turned a parallel combination into one addition.

\section{Phasor Lead and Lag}
In a capacitor, current leads voltage by \(90^\circ\) (because
\(i=C\,\dd v/\dd t\) is proportional to the derivative); in an inductor, voltage
leads current by \(90^\circ\). A positive impedance angle means net inductive
behavior; a negative angle means net capacitive behavior.

\begin{exambox}
If a question says ``lead'' or ``lag,'' sketch a tiny phasor diagram before
answering. Most mistakes come from recalling the right \(90^\circ\) relationship
with the direction reversed. Mnemonic: ELI the ICE man---in an inductor (L)
voltage E leads current I; in a capacitor (C) current I leads voltage E.
\end{exambox}

\section{Complex Power}
With RMS phasors and the passive sign convention, complex power uses the
conjugate of the current:
\[
S=V I^{*}=P+\jj Q .
\]
The real part \(P\) is average power actually dissipated (watts); the imaginary
part \(Q\) is reactive power (VAR), the energy sloshing back and forth between the
source and the reactive elements each cycle. The conjugate is what makes \(P\)
come out as the in-phase product.

\subsection{Worked Example: Real and Reactive Power}
A load draws \(I=2\angle{-30^\circ}\ \si{\ampere}\) (RMS) from
\(V=120\angle0^\circ\ \si{\volt}\) (RMS)---current lagging, so the load is
inductive. Then
\[
S=VI^{*}=120\angle0^\circ\cdot 2\angle{+30^\circ}=240\angle30^\circ
=207.8+\jj120\ \mathrm{VA}.
\]
So \(P=\SI{207.8}{\watt}\) is dissipated and \(Q=120\,\mathrm{VAR}\) is exchanged with
the load's reactance. A purely reactive load (\(\pm90^\circ\)) would show
\(P=0\): it consumes no average power.

\section{Worked Example: Series Reactance}
Find the impedance of a \(\SI{50}{\ohm}\) resistor in series with
\(\SI{10}{\micro\henry}\) at \(\SI{7.1}{\mega\hertz}\):
\[
\omega=2\pi(7.1\times10^6)=\SI{4.46e7}{\radian\per\second},
\qquad
X_L=\omega L\approx\SI{446}{\ohm}.
\]
Thus \(Z=50+\jj446\ \Omega\), with \(|Z|=\SI{449}{\ohm}\) and
\(\theta\approx83.6^\circ\)---strongly inductive, since the reactance dwarfs the
resistance.
